\chapter{Couche limite laminaire : la méthode de Thwaites}\label{ch:thwaites}
\iterbanner{V0}

\section{De von Kármán à une équation fermée}
On multiplie l'équation de von Kármán par $2U_e\theta/\nu$ :
\[
\frac{U_e}{\nu}\frac{\dd\theta^2}{\dd s}+2\big(2+H\big)\frac{\theta^2}{\nu}\frac{\dd U_e}{\dd s}=\frac{U_e\,\theta\,c_f}{\nu}.
\]
Deux groupements sans dimension apparaissent. Le premier est le \textbf{paramètre de gradient de pression}
\[
\lambda=\frac{\theta^2}{\nu}\frac{\dd U_e}{\dd s}.
\]
Le second est le frottement réduit $\ell=\tfrac12c_f\,\Re_\theta=\tfrac{\tau_w\theta}{\mu U_e}$. L'équation devient
\[
\frac{U_e}{\nu}\frac{\dd\theta^2}{\dd s}=2\big[\ell-(2+H)\lambda\big]\equiv F .
\]

\paragraph{L'hypothèse de Thwaites.} En laminaire, la forme du profil de vitesse est gouvernée principalement par le gradient de pression local. On suppose donc que $H$ et $\ell$ sont des fonctions du seul $\lambda$ : c'est une famille de profils à un paramètre. Thwaites a tracé $F(\lambda)$ pour toutes les solutions exactes connues (Falkner-Skan, Howarth, etc.) et constaté que les points s'alignent presque sur une droite (figure~\ref{fig:thwaites-f}) :
\[
F(\lambda)\approx a-b\lambda,\qquad a=0{,}45,\quad b=6 .
\]

\paragraph{La quadrature.} Avec cette droite, l'équation devient linéaire en $\theta^2$ :
\[
U_e\frac{\dd\theta^2}{\dd s}+6\,\theta^2\frac{\dd U_e}{\dd s}=0{,}45\,\nu
\quad\Longleftrightarrow\quad
\frac{1}{U_e^5}\frac{\dd}{\dd s}\big(\theta^2U_e^6\big)=0{,}45\,\nu .
\]
On a multiplié par $U_e^5$ pour faire apparaître une dérivée exacte. C'est l'origine de la puissance 6, et donc du 5 dans l'intégrale. En partant du point d'arrêt, où $U_e=0$ :
\[
\boxed{\ \theta^2(s)=\frac{0{,}45\,\nu}{U_e^6(s)}\int_0^sU_e^5(s')\,\dd s'\ }
\]
On n'a plus d'EDO à résoudre : $\theta$ s'obtient par une simple intégrale cumulée de $U_e^5$.

\section{Les corrélations de fermeture}
Une fois $\theta$ connu, on calcule $\lambda$, puis $H$ et $c_f$ par les corrélations de Thwaites, sous la forme de Cebeci et Bradshaw :
\[
0\le\lambda\le0{,}1:\quad \ell=0{,}22+1{,}57\lambda-1{,}8\lambda^2,\qquad H=2{,}61-3{,}75\lambda+5{,}24\lambda^2,
\]
\[
-0{,}1\le\lambda<0:\quad \ell=0{,}22+1{,}402\lambda+\frac{0{,}018\lambda}{\lambda+0{,}107},\qquad H=2{,}088+\frac{0{,}0731}{\lambda+0{,}14},
\]
\[
c_f=\frac{2\,\ell(\lambda)}{\Re_\theta},\qquad \Re_\theta=\frac{U_e\theta}{\nu}.
\]
Le \textbf{décollement laminaire} correspond à $\ell=0$, soit $\lambda\approx-0{,}09$ (constante \code{LAMBDA\_SEP}).

\begin{figure}[htbp]
\centering
% figure générée par tools/make_figures.py : thwaites_fermeture
\begin{tikzpicture}
\begin{groupplot}[group style={group size=2 by 1,horizontal sep=1.6cm},width=0.5\linewidth,height=5.4cm,
  xlabel={$\lambda=\theta^2U_e'/\nu$},xmin=-0.1,xmax=0.1,title style={font=\small},tick label style={font=\scriptsize},
  legend style={draw=none,font=\scriptsize}]
\nextgroupplot[title={Corrélations de fermeture},legend pos=north east,ymin=0,ymax=4]
\addplot[cblue,very thick] coordinates {(-0.09,3.55) (-0.088403,3.5048) (-0.086807,3.4622) (-0.08521,3.4222) (-0.083613,3.3844) (-0.082017,3.3487) (-0.08042,3.3149) (-0.078824,3.2829) (-0.077227,3.2525) (-0.07563,3.2236) (-0.074034,3.1961) (-0.072437,3.17) (-0.07084,3.145) (-0.069244,3.1211) (-0.067647,3.0983) (-0.06605,3.0765) (-0.064454,3.0556) (-0.062857,3.0356) (-0.061261,3.0164) (-0.059664,2.9979) (-0.058067,2.9802) (-0.056471,2.9631) (-0.054874,2.9467) (-0.053277,2.9309) (-0.051681,2.9157) (-0.050084,2.901) (-0.048487,2.8868) (-0.046891,2.8731) (-0.045294,2.8599) (-0.043697,2.8471) (-0.042101,2.8347) (-0.040504,2.8227) (-0.038908,2.8111) (-0.037311,2.7999) (-0.035714,2.789) (-0.034118,2.7784) (-0.032521,2.7681) (-0.030924,2.7582) (-0.029328,2.7485) (-0.027731,2.7391) (-0.026134,2.73) (-0.024538,2.7211) (-0.022941,2.7125) (-0.021345,2.7041) (-0.019748,2.6959) (-0.018151,2.6879) (-0.016555,2.6802) (-0.014958,2.6726) (-0.013361,2.6652) (-0.011765,2.658) (-0.010168,2.651) (-0.0085714,2.6442) (-0.0069748,2.6375) (-0.0053782,2.631) (-0.0037815,2.6246) (-0.0021849,2.6184) (-0.00058824,2.6123) (0.0010084,2.6062) (0.002605,2.6003) (0.0042017,2.5943) (0.0057983,2.5884) (0.007395,2.5826) (0.0089916,2.5767) (0.010588,2.5709) (0.012185,2.5651) (0.013782,2.5593) (0.015378,2.5536) (0.016975,2.5479) (0.018571,2.5422) (0.020168,2.5365) (0.021765,2.5309) (0.023361,2.5253) (0.024958,2.5197) (0.026555,2.5141) (0.028151,2.5086) (0.029748,2.5031) (0.031345,2.4976) (0.032941,2.4922) (0.034538,2.4867) (0.036134,2.4813) (0.037731,2.476) (0.039328,2.4706) (0.040924,2.4653) (0.042521,2.46) (0.044118,2.4548) (0.045714,2.4495) (0.047311,2.4443) (0.048908,2.4391) (0.050504,2.434) (0.052101,2.4288) (0.053697,2.4237) (0.055294,2.4187) (0.056891,2.4136) (0.058487,2.4086) (0.060084,2.4036) (0.061681,2.3986) (0.063277,2.3937) (0.064874,2.3888) (0.066471,2.3839) (0.068067,2.379) (0.069664,2.3742) (0.071261,2.3694) (0.072857,2.3646) (0.074454,2.3598) (0.07605,2.3551) (0.077647,2.3504) (0.079244,2.3457) (0.08084,2.3411) (0.082437,2.3365) (0.084034,2.3319) (0.08563,2.3273) (0.087227,2.3228) (0.088824,2.3183) (0.09042,2.3138) (0.092017,2.3093) (0.093613,2.3049) (0.09521,2.3005) (0.096807,2.2961) (0.098403,2.2917) (0.1,2.2874)};\addlegendentry{$H(\lambda)$}
\addplot[corange,very thick] coordinates {(-0.09,-0.014741) (-0.088403,0.10491) (-0.086807,0.20919) (-0.08521,0.30146) (-0.083613,0.38419) (-0.082017,0.45921) (-0.08042,0.5279) (-0.078824,0.59135) (-0.077227,0.65039) (-0.07563,0.7057) (-0.074034,0.75782) (-0.072437,0.80719) (-0.07084,0.85418) (-0.069244,0.89909) (-0.067647,0.94217) (-0.06605,0.98364) (-0.064454,1.0237) (-0.062857,1.0624) (-0.061261,1.1) (-0.059664,1.1366) (-0.058067,1.1723) (-0.056471,1.2071) (-0.054874,1.2412) (-0.053277,1.2745) (-0.051681,1.3073) (-0.050084,1.3394) (-0.048487,1.371) (-0.046891,1.4022) (-0.045294,1.4329) (-0.043697,1.4631) (-0.042101,1.493) (-0.040504,1.5225) (-0.038908,1.5517) (-0.037311,1.5805) (-0.035714,1.6091) (-0.034118,1.6374) (-0.032521,1.6655) (-0.030924,1.6933) (-0.029328,1.7209) (-0.027731,1.7482) (-0.026134,1.7754) (-0.024538,1.8024) (-0.022941,1.8292) (-0.021345,1.8559) (-0.019748,1.8824) (-0.018151,1.9087) (-0.016555,1.935) (-0.014958,1.961) (-0.013361,1.987) (-0.011765,2.0128) (-0.010168,2.0385) (-0.0085714,2.0642) (-0.0069748,2.0897) (-0.0053782,2.1151) (-0.0037815,2.1404) (-0.0021849,2.1656) (-0.00058824,2.1908) (0.0010084,2.2158) (0.002605,2.2408) (0.0042017,2.2656) (0.0057983,2.2904) (0.007395,2.3151) (0.0089916,2.3397) (0.010588,2.3642) (0.012185,2.3886) (0.013782,2.413) (0.015378,2.4372) (0.016975,2.4613) (0.018571,2.4854) (0.020168,2.5093) (0.021765,2.5332) (0.023361,2.5569) (0.024958,2.5806) (0.026555,2.6042) (0.028151,2.6277) (0.029748,2.6511) (0.031345,2.6744) (0.032941,2.6976) (0.034538,2.7208) (0.036134,2.7438) (0.037731,2.7668) (0.039328,2.7896) (0.040924,2.8124) (0.042521,2.835) (0.044118,2.8576) (0.045714,2.8801) (0.047311,2.9025) (0.048908,2.9248) (0.050504,2.947) (0.052101,2.9691) (0.053697,2.9911) (0.055294,3.0131) (0.056891,3.0349) (0.058487,3.0567) (0.060084,3.0783) (0.061681,3.0999) (0.063277,3.1214) (0.064874,3.1428) (0.066471,3.1641) (0.068067,3.1853) (0.069664,3.2064) (0.071261,3.2274) (0.072857,3.2483) (0.074454,3.2691) (0.07605,3.2899) (0.077647,3.3105) (0.079244,3.3311) (0.08084,3.3516) (0.082437,3.3719) (0.084034,3.3922) (0.08563,3.4124) (0.087227,3.4325) (0.088824,3.4525) (0.09042,3.4724) (0.092017,3.4923) (0.093613,3.512) (0.09521,3.5316) (0.096807,3.5512) (0.098403,3.5706) (0.1,3.59)};\addlegendentry{$10\,\ell(\lambda)$}
\addplot[cmuted,dashed] coordinates {(-0.09,0) (-0.09,4)};
\node[cmuted,font=\scriptsize,rotate=90,anchor=south] at (axis cs:-0.09,2.6) {décollement};
\nextgroupplot[title={$F(\lambda)=2[\ell-(2+H)\lambda]$ et sa droite},legend pos=north east]
\addplot[cblue,very thick] coordinates {(-0.09,0.99605) (-0.088403,0.99426) (-0.086807,0.99015) (-0.08521,0.98434) (-0.083613,0.97726) (-0.082017,0.96921) (-0.08042,0.96043) (-0.078824,0.9511) (-0.077227,0.94135) (-0.07563,0.93127) (-0.074034,0.92094) (-0.072437,0.91043) (-0.07084,0.89978) (-0.069244,0.88903) (-0.067647,0.87821) (-0.06605,0.86734) (-0.064454,0.85644) (-0.062857,0.84553) (-0.061261,0.83462) (-0.059664,0.82372) (-0.058067,0.81283) (-0.056471,0.80197) (-0.054874,0.79113) (-0.053277,0.78032) (-0.051681,0.76955) (-0.050084,0.75881) (-0.048487,0.74811) (-0.046891,0.73744) (-0.045294,0.72682) (-0.043697,0.71623) (-0.042101,0.70568) (-0.040504,0.69518) (-0.038908,0.68471) (-0.037311,0.67428) (-0.035714,0.66389) (-0.034118,0.65354) (-0.032521,0.64322) (-0.030924,0.63294) (-0.029328,0.6227) (-0.027731,0.61249) (-0.026134,0.60232) (-0.024538,0.59217) (-0.022941,0.58207) (-0.021345,0.57199) (-0.019748,0.56195) (-0.018151,0.55193) (-0.016555,0.54195) (-0.014958,0.53199) (-0.013361,0.52207) (-0.011765,0.51217) (-0.010168,0.50229) (-0.0085714,0.49245) (-0.0069748,0.48262) (-0.0053782,0.47283) (-0.0037815,0.46305) (-0.0021849,0.4533) (-0.00058824,0.44358) (0.0010084,0.43387) (0.002605,0.42419) (0.0042017,0.41452) (0.0057983,0.40488) (0.007395,0.39525) (0.0089916,0.38564) (0.010588,0.37605) (0.012185,0.36648) (0.013782,0.35692) (0.015378,0.34739) (0.016975,0.33787) (0.018571,0.32836) (0.020168,0.31888) (0.021765,0.30941) (0.023361,0.29996) (0.024958,0.29052) (0.026555,0.2811) (0.028151,0.2717) (0.029748,0.26231) (0.031345,0.25293) (0.032941,0.24358) (0.034538,0.23423) (0.036134,0.2249) (0.037731,0.21558) (0.039328,0.20628) (0.040924,0.19699) (0.042521,0.18772) (0.044118,0.17846) (0.045714,0.16921) (0.047311,0.15997) (0.048908,0.15074) (0.050504,0.14153) (0.052101,0.13233) (0.053697,0.12314) (0.055294,0.11396) (0.056891,0.1048) (0.058487,0.095641) (0.060084,0.086495) (0.061681,0.07736) (0.063277,0.068234) (0.064874,0.059119) (0.066471,0.050012) (0.068067,0.040916) (0.069664,0.031828) (0.071261,0.022748) (0.072857,0.013677) (0.074454,0.0046148) (0.07605,-0.0044399) (0.077647,-0.013487) (0.079244,-0.022526) (0.08084,-0.031559) (0.082437,-0.040584) (0.084034,-0.049603) (0.08563,-0.058615) (0.087227,-0.067621) (0.088824,-0.076622) (0.09042,-0.085616) (0.092017,-0.094606) (0.093613,-0.10359) (0.09521,-0.11257) (0.096807,-0.12154) (0.098403,-0.13051) (0.1,-0.13948)};\addlegendentry{$F$ exact (corrélations)}
\addplot[corange,dashed,very thick,domain=-0.09:0.1] {0.45-6*x};\addlegendentry{$0{,}45-6\lambda$}
\end{groupplot}
\end{tikzpicture}

\caption{À gauche, les corrélations $H(\lambda)$ et $\ell(\lambda)$ utilisées par le code. À droite, la fonction $F(\lambda)$ qu'elles impliquent, comparée à la droite $0{,}45-6\lambda$ de Thwaites. L'écart reste faible sur toute la plage utile, ce qui justifie la quadrature.}\label{fig:thwaites-f}
\end{figure}

\paragraph{Intuition physique de $\lambda$.} $\theta^2/\nu$ est le temps de diffusion visqueuse à travers la couche, et $1/U_e'$ le temps caractéristique de variation de l'écoulement extérieur. $\lambda$ compare les deux :
\begin{itemize}
\item $\lambda>0$ (accélération) : profil plein, $H$ petit, frottement élevé ;
\item $\lambda<0$ (décélération) : profil creusé près de la paroi, $H$ grand ;
\item $\lambda\approx-0{,}09$ : le frottement s'annule, la couche décolle.
\end{itemize}

\section{Deux cas exacts}
\paragraph{Plaque plane ($U_e=1$).} L'intégrale vaut $s$, d'où $\theta^2=0{,}45\,\nu s$, soit
\[
\theta=0{,}671\,\frac{s}{\sqrt{\Re_s}},
\]
contre $0{,}664$ pour la solution exacte de Blasius : 1~\% d'écart. Avec $\lambda=0$, on obtient $H=2{,}61$ (Blasius : 2,59).

\paragraph{Point d'arrêt ($U_e=ks$, Hiemenz).} L'intégrale vaut $k^5s^6/6$, d'où $\theta^2=0{,}075\,\nu/k$ : l'épaisseur est \emph{constante} près du point d'arrêt, et $\lambda=0{,}075$. C'est ce qui permet de démarrer la quadrature en $s=0$ sans singularité.

\begin{figure}[htbp]
\centering
% figure générée par tools/make_figures.py : plaque_plane
\begin{tikzpicture}
\begin{groupplot}[group style={group size=2 by 1,horizontal sep=1.6cm},width=0.5\linewidth,height=5.2cm,
  xlabel={$x/c$},xmin=0,xmax=1,title style={font=\small},tick label style={font=\scriptsize},
  legend style={draw=none,font=\scriptsize},legend pos=north west]
\nextgroupplot[title={Laminaire (Thwaites), $Re=10^7$},ylabel={$10^3\,\theta/c$}]
\addplot[cblue,very thick] coordinates {(0.0001,0.0015) (0.020148,0.030074) (0.040196,0.042504) (0.060244,0.052046) (0.080292,0.060091) (0.10034,0.067179) (0.12039,0.073588) (0.14044,0.079482) (0.16048,0.084968) (0.18053,0.090121) (0.20058,0.094994) (0.22063,0.09963) (0.24068,0.10406) (0.26073,0.10831) (0.28077,0.11239) (0.30082,0.11634) (0.32087,0.12015) (0.34092,0.12385) (0.36097,0.12744) (0.38101,0.13093) (0.40106,0.13433) (0.42111,0.13765) (0.44116,0.14089) (0.46121,0.14406) (0.48125,0.14715) (0.5013,0.15019) (0.52135,0.15316) (0.5414,0.15608) (0.56145,0.15894) (0.5815,0.16176) (0.60154,0.16452) (0.62159,0.16724) (0.64164,0.16992) (0.66169,0.17255) (0.68174,0.17515) (0.70178,0.1777) (0.72183,0.18022) (0.74188,0.18271) (0.76193,0.18516) (0.78198,0.18758) (0.80202,0.18997) (0.82207,0.19233) (0.84212,0.19466) (0.86217,0.19697) (0.88222,0.19924) (0.90227,0.20149) (0.92231,0.20372) (0.94236,0.20592) (0.96241,0.2081) (0.98246,0.21026)};\addlegendentry{Thwaites (code)}
\addplot[cgreen,dashed,very thick] coordinates {(0.0001,0.0020998) (0.020148,0.029805) (0.040196,0.042098) (0.060244,0.051538) (0.080292,0.059498) (0.10034,0.066513) (0.12039,0.072855) (0.14044,0.078688) (0.16048,0.084117) (0.18053,0.089217) (0.20058,0.09404) (0.22063,0.098628) (0.24068,0.10301) (0.26073,0.10722) (0.28077,0.11126) (0.30082,0.11517) (0.32087,0.11894) (0.34092,0.1226) (0.36097,0.12615) (0.38101,0.12961) (0.40106,0.13298) (0.42111,0.13626) (0.44116,0.13947) (0.46121,0.1426) (0.48125,0.14567) (0.5013,0.14867) (0.52135,0.15161) (0.5414,0.1545) (0.56145,0.15733) (0.5815,0.16012) (0.60154,0.16286) (0.62159,0.16555) (0.64164,0.1682) (0.66169,0.1708) (0.68174,0.17337) (0.70178,0.1759) (0.72183,0.1784) (0.74188,0.18086) (0.76193,0.18328) (0.78198,0.18568) (0.80202,0.18805) (0.82207,0.19038) (0.84212,0.19269) (0.86217,0.19497) (0.88222,0.19722) (0.90227,0.19945) (0.92231,0.20165) (0.94236,0.20383) (0.96241,0.20599) (0.98246,0.20813)};\addlegendentry{Blasius}
\nextgroupplot[title={Turbulent dès $x=0$ (Head)}]
\addplot[corange,very thick] coordinates {(0.0001,0.0015) (0.020148,0.054417) (0.040196,0.097004) (0.060244,0.13614) (0.080292,0.17316) (0.10034,0.20865) (0.12039,0.24297) (0.14044,0.27633) (0.16048,0.3089) (0.18053,0.34077) (0.20058,0.37205) (0.22063,0.4028) (0.24068,0.43307) (0.26073,0.4629) (0.28077,0.49235) (0.30082,0.52143) (0.32087,0.55017) (0.34092,0.57861) (0.36097,0.60676) (0.38101,0.63464) (0.40106,0.66226) (0.42111,0.68964) (0.44116,0.71679) (0.46121,0.74373) (0.48125,0.77047) (0.5013,0.79701) (0.52135,0.82336) (0.5414,0.84954) (0.56145,0.87554) (0.5815,0.90139) (0.60154,0.92707) (0.62159,0.95261) (0.64164,0.978) (0.66169,1.0033) (0.68174,1.0284) (0.70178,1.0534) (0.72183,1.0782) (0.74188,1.1029) (0.76193,1.1276) (0.78198,1.1521) (0.80202,1.1765) (0.82207,1.2008) (0.84212,1.2249) (0.86217,1.249) (0.88222,1.273) (0.90227,1.2969) (0.92231,1.3207) (0.94236,1.3444) (0.96241,1.368) (0.98246,1.3915)};\addlegendentry{Head (code)}
\addplot[cgreen,dashed,very thick] coordinates {(0.0001,0.00090428) (0.020148,0.063051) (0.040196,0.10956) (0.060244,0.15144) (0.080292,0.19056) (0.10034,0.22776) (0.12039,0.26349) (0.14044,0.29805) (0.16048,0.33163) (0.18053,0.36437) (0.20058,0.3964) (0.22063,0.42779) (0.24068,0.45862) (0.26073,0.48893) (0.28077,0.51879) (0.30082,0.54821) (0.32087,0.57725) (0.34092,0.60593) (0.36097,0.63427) (0.38101,0.6623) (0.40106,0.69004) (0.42111,0.7175) (0.44116,0.74469) (0.46121,0.77165) (0.48125,0.79837) (0.5013,0.82486) (0.52135,0.85115) (0.5414,0.87724) (0.56145,0.90313) (0.5815,0.92884) (0.60154,0.95437) (0.62159,0.97973) (0.64164,1.0049) (0.66169,1.03) (0.68174,1.0549) (0.70178,1.0796) (0.72183,1.1042) (0.74188,1.1287) (0.76193,1.153) (0.78198,1.1772) (0.80202,1.2013) (0.82207,1.2253) (0.84212,1.2491) (0.86217,1.2728) (0.88222,1.2965) (0.90227,1.32) (0.92231,1.3434) (0.94236,1.3667) (0.96241,1.3899) (0.98246,1.413)};\addlegendentry{loi en $1/7$}
\end{groupplot}
\end{tikzpicture}

\caption{Validation sur plaque plane à $\Re=10^7$ (tests V2 et V3 de \code{validation.py}). À gauche, Thwaites contre Blasius : 1,0~\% d'écart, la différence entre les constantes 0,671 et 0,664. À droite, Head contre la loi turbulente en $1/7$ (ch.~\ref{ch:head}) : moins de 5~\% d'écart au-delà de $x=0{,}3$.}\label{fig:plaque}
\end{figure}

\begin{codebox}[ : bloc Thwaites de \code{march\_surface}]
\begin{lstlisting}
s0 = np.concatenate([[0.0], s])          # on ajoute le point d'arret, s = 0
U0 = np.concatenate([[0.0], Ue])         # ou Ue = 0
integ = np.concatenate([[0.0], np.cumsum(0.5 * (U0[1:]**5 + U0[:-1]**5) * np.diff(s0))])[1:]
theta_lam = np.sqrt(0.45 / Re * integ / Ue**6)       # nu = 1/Re
lam = theta_lam**2 * dUe * Re                         # dUe = np.gradient(Ue, s)
H_lam, ell = thwaites_closure(lam)
Re_th = Ue * theta_lam * Re
\end{lstlisting}
La quadrature est une méthode des trapèzes sur $U_e^5$. Le point $s=0$, $U_e=0$ est ajouté explicitement : sans lui, l'intégrale démarrerait à la première collocation et $\theta$ serait faux sur les premières stations. \code{thwaites\_closure} borne $\lambda$ à $[-0{,}1;\,0{,}1]$ avant d'évaluer les corrélations.
\end{codebox}

\begin{piege}
$\theta$ dépend de tout l'historique amont (l'intégrale), alors que $\lambda$, donc $H$, est \emph{local}. Thwaites réagit donc mal aux variations brutales de $U_e$ : bulbe de décollement, choc, ou la chute de vitesse au bord de fuite fermé. De plus, au-delà de $\lambda<-0{,}09$, la méthode n'a plus de sens. On ne la prolonge pas : on bascule en turbulent (ch.~\ref{ch:transition}).
\end{piege}

\exercices
\begin{exercice}
Retrouver le résultat de Hiemenz : pour $U_e=ks$, calculer $\theta^2$ par la formule de Thwaites, puis $\lambda$. Que vaut $H$ au point d'arrêt d'après les corrélations ?
\end{exercice}
\begin{exercice}
Pourquoi la puissance 5 ? Reprendre le calcul avec une droite quelconque $F=a-b\lambda$ et montrer qu'on obtient $\theta^2U_e^{b}=a\,\nu\int_0^sU_e^{b-1}\,\dd s'$.
\end{exercice}
